\documentclass[12pt, a4paper, oneside]{ctexart}
\usepackage{amsmath, amsthm, amssymb, bm, xcolor, tikz, framed, graphicx, hyperref, mathrsfs}
\usepackage{bookmark}
\usepackage{pgfplots}
\pgfplotsset{compat=1.18}

\newcommand{\bra}[1]{\langle #1 |}
\newcommand{\ket}[1]{| #1 \rangle}
\newcommand{\braket}[2]{\langle #1 | #2 \rangle}

\title{Quantum Mechanics Homework (Week 17)}
\author{Tang Yi \\ [学号已脱敏]}
\usepackage[a4paper, left=2cm, right=2cm, top=2cm, bottom=2cm]{geometry}
\linespread{1.3} % 行距更紧凑
\definecolor{shadecolor}{RGB}{241, 241, 255}

\newcounter{problemname}
\newenvironment{problem}
  {\begin{shaded}\stepcounter{problemname}\par\noindent\textbf{题目 \theproblemname. }}
  {\end{shaded}\par}
\newenvironment{solution}
  {\par\noindent\textbf{题目 \theproblemname 的解答. }}
  {\par}
\newenvironment{note}
  {\par\noindent\textbf{题目 \theproblemname 的注记. }}
  {\par}

\begin{document}

\maketitle


\begin{problem}
  在 w 16-1 中，若考虑 $x$ 方向的磁场是缓慢加上力的：

$$B_x(t) = B_1 e^{\frac{t}{\tau}} (-\infty < t < 0)$$

其中将初始时刻取为 $t = -\infty$，求 $t = 0$ 时粒子处于沿磁场方向自旋向上态的概率，并说明是否满足绝热定理。
\end{problem}

\begin{solution}
仍考虑自旋 $1/2$ 粒子，其哈密顿量为
\[
H(t)=\mu_B B_0\sigma_z+\mu_B B_x(t)\sigma_x,
\qquad 
B_x(t)=B_1 e^{t/\tau}\;(t<0),
\]
其中 $B_1\ll B_0$，并取初始时刻 $t=-\infty$。
此时
\[
B_x(-\infty)=0,
\]
哈密顿量趋于
\[
H(-\infty)=\mu_B B_0\sigma_z.
\]
初态取为
\[
|\psi(-\infty)\rangle=|\downarrow_z\rangle,
\]
即沿 $z$ 方向磁场的低能态。

在任意时刻 $t$，哈密顿量可以写成
\[
H(t)=\mu_B \mathbf B(t)\cdot\boldsymbol{\sigma},
\qquad 
\mathbf B(t)=(B_x(t),0,B_0).
\]
其瞬时本征态是沿瞬时磁场方向
\[
\hat{\mathbf n}(t)
=\frac{1}{\sqrt{B_0^2+B_x^2(t)}}(B_x(t),0,B_0)
\]
的自旋上、下态。对应的能级差为
\[
\Delta E(t)=2\mu_B\sqrt{B_0^2+B_x^2(t)}.
\]

绝热定理要求
\[
\left|
\frac{\langle +(t)|\dot H(t)|-(t)\rangle}
{[\Delta E(t)]^2}
\right|
\ll 1.
\]
在本问题中
\[
\dot B_x(t)=\frac{1}{\tau}B_x(t),
\]
可估算绝热条件为
\[
\frac{\mu_B |\dot B_x(t)|}{(2\mu_B B_0)^2}
\sim
\frac{B_x(t)}{4B_0^2\tau}
\ll 1.
\]
在整个演化过程中 $B_x(t)\le B_1\ll B_0$，因此若
\[
\tau \gg \frac{1}{\omega_0},
\qquad
\omega_0=\frac{2\mu_B B_0}{\hbar},
\]
体系满足绝热条件。

在绝热近似成立时，体系始终处于瞬时本征态。
由于初态 $|\downarrow_z\rangle$ 在 $t=-\infty$ 正是瞬时磁场方向的自旋下态，
演化到 $t=0$ 时，体系将演化为沿最终磁场
\[
\mathbf B(0)=(B_1,0,B_0)
\]
方向的自旋下态（仅相差一个相位）。

因此，在 $t=0$ 时粒子处于沿磁场方向自旋向上态的概率为
\[
P_{\uparrow}^{(\text{adiabatic})}(0)=0.
\]

若考虑非绝热修正，在 $B_1\ll B_0$、$\tau$ 足够大时，跃迁概率是指数小量，
可估计为
\[
P_{\uparrow}(0)\sim 
\mathcal O\!\left(\frac{1}{(\omega_0\tau)^2}\right)\ll1.
\]

结论是：当磁场在时间尺度 $\tau\gg 1/\omega_0$ 上缓慢加上时，
体系满足绝热定理，粒子始终跟随瞬时本征态演化，
在 $t=0$ 时几乎不可能出现在沿磁场方向的自旋向上态。

\end{solution}

\end{document}